% Generated draft: requires mathematical and source review.
\documentclass[11pt]{article}
\usepackage[margin=25mm]{geometry}
\usepackage{fontspec}
\setmainfont{Latin Modern Roman}
\usepackage{amsmath,amssymb,mathtools,unicode-math}
\setmathfont{Latin Modern Math}
\usepackage{enumitem,needspace,xurl,hyperref}
\setlength{\parindent}{0pt}
\setlength{\parskip}{5pt}
\setlength{\emergencystretch}{3em}
\newcommand{\N}{\mathbb{N}}
\newcommand{\Z}{\mathbb{Z}}
\newcommand{\Q}{\mathbb{Q}}
\newcommand{\R}{\mathbb{R}}
\newcommand{\C}{\mathbb{C}}
\newcommand{\F}{\mathbb{F}}
\newcommand{\A}{\mathbb{A}}
\newcommand{\PP}{\mathbb{P}}
\newcommand{\E}{\mathbb{E}}
\newcommand{\eps}{\varepsilon}
\newcommand{\dd}{\,\mathrm{d}}
\begin{document}
\section*{UnsolvedMath ID: 30000974}
\textbf{Source problem code:} catalogue entry
\subsection{Definitions and mathematical statement}
Let \(X\) be a smooth projective toric variety over the intended field \(k\), let \(L\) be an ample line bundle and let \(L'\) be nef. Here nef means that the degree on every integral projective curve is nonnegative. The source asks:
\begin{enumerate}
\item For every such \((X,L,L')\), is the multiplication-of-sections map
\[
H^0(X,L)\otimes_k H^0(X,L')\longrightarrow H^0(X,L\otimes L'),
\qquad s\otimes t\longmapsto st,
\]
surjective? This is the supplied Oda problem.
\item With \(X\) smooth and \(L\) ample, if the section ring \(R(X,L)\) is normal, is its ideal of relations generated in degree two? This is the supplied Sturmfels question.
\end{enumerate}
To make the second question precise one must fix the ring convention. The complete section ring is
\[
R(X,L)=\bigoplus_{m\geq0}H^0(X,L^{\otimes m}).
\]
There is a natural graded map \(\operatorname{Sym}_k H^0(X,L)\to R(X,L)\); its kernel is the ideal of relations among degree-one sections. When this map is surjective it is a presentation of the complete section ring, and quadratic generation means that its kernel is generated by its degree-two part. The degree-one-generated homogeneous coordinate ring is the image of this map. Normality, meaning integral closedness of the chosen ring in its fraction field, must be distinguished from normal generation of \(L\), meaning surjectivity of all degree-one multiplication maps. The source does not specify which convention its normality assumption and relation ideal use; neither an extra generation hypothesis nor a replacement of its stated hypothesis is silently made.
\subsection{Short English statement}
Ask whether ample-by-nef multiplication is always surjective on a smooth projective toric variety, and whether the stated normality hypothesis forces quadratic defining relations. The second question requires care about the difference between the complete section ring, its degree-one-generated coordinate ring, and normal generation.
\subsection{Sources}
\begin{itemize}
\item Local UnsolvedMath registry, permanent ID \(30000974\), source code \texttt{catalogue entry}, ``Oda and Sturmfels Conjectures for Smooth Toric Varieties.'' Registry curation: Ulam AI / UnsolvedMath Contributors, CC BY 4.0; underlying sources retain their supplied provenance.
\item \emph{Classical Algebraic Geometry} (2009), Oberwolfach Reports, pp.\ 1425--1474, DOI \url{https://doi.org/10.4171/owr/2008/26}.
\item Sturmfels, supplied title \emph{Equations defining toric varieties}: \url{https://arxiv.org/abs/alg-geom/9610018}.
\item The supplied assessment checked 2026-08-21 labels the questions open in general for smooth toric varieties, mentioning Oda normal generation, Sturmfels quadratic relations, Bøgvad normal-Koszul predictions, and diagonal-splitting and low-dimensional cases. The Bøgvad prediction is background in this record, not an extra requested conjecture; the partial results and status are dated supplied claims.
\end{itemize}
\end{document}
